To begin, we analyze this hiring model 
for a single group of candidates.
The employer's goal is to minimize the expected loss, $l_{\alpha}(\pi)$, 
while maintaining a given acceptance probability.
For brevity, we relegate all proofs to the appendix.

\subsection{The Simple Threshold Policy (Equal Number of Tests)}

Consider the setting where 
the employer must subject all candidates 
to an equal number of tests $\tau$ and threshold $\theta$ 
(these parameters are chosen by the employer but thereafter constant across candidates).
For a given threshold, we can relate
the flip probability (error rate) of the test
to the probability that a candidate is accepted as follows:

Recall that $\hat{y}_{i,j}=y_i\xor Br(\eta)$, $S_\tau=\sum_{j=1}^\tau \hat{y}_{i,j}$, that $Z^\eta_\tau=\sum_{t=1}^\tau \mathbb{I}(\hat{y}_{i,j}\neq y_i)$, and that $\tau$ and $\theta$ are the only parameters of the threshold policy, $\pi$.
Informally, $S_\tau$ is the number of passed tests
and $Z^\eta_\tau$ is the number of flips 
(tests in error).
The probability of hiring an unskilled candidate is given by:
$$\Pr[\pi(\hat{y}_{i,1},\dots,\hat{y}_{i,\tau})=1|y_i=0]=\Pr[S_{\tau}\geq \theta|y_i=0]=
\Pr[Z_{\tau}^{\eta}\geq \theta].$$

Since $Z_{\tau}^{\eta}$ is a binomial random variable with parameters $\tau$ and $\eta$, we can calculate this probability precisely as:
$\Pr[\pi(\hat{y}_{i,1},\dots,\hat{y}_{i,\tau})=1|y_i=0]=\Pr[Z_{\tau}^{\eta}\geq \theta]= \sum_{k=\theta}^\tau{\tau \choose k}\eta^k(1-\eta)^{\tau-k}$,
and the probability of rejecting a skilled candidate is the probability that they encounter more than $\tau-\theta$ flips, thus: 
$\Pr[\pi(\hat{y}_{i,1},\dots,\hat{y}_{i,\tau})=0|y_i=+1]=\Pr[S_{\tau}<\theta|y=+1]=\Pr[Z_\tau^{\eta}> \tau- \theta]=$\\
$\sum_{k=\tau- \theta+1}^{\tau}{\tau \choose k}\eta^k(1-\eta)^{\tau-k}.$
Similarly, given a candidate's skill level, 
we can calculate the probability 
that 
they obtain exactly $k$ positive tests 
out of $\tau$, i.e,
$$\Pr[S_{\tau}=k|y_i=0]=\Pr[Z_\tau^{\eta}= k]={\tau \choose k}\eta^{k}(1-\eta)^{\tau-k}.$$
$$\Pr[S_{\tau}=k|y_i=+1]=\Pr[Z_\tau^{\eta}= \tau- k]={\tau \choose k}\eta^{\tau-k}(1-\eta)^{k}.$$
% \end{observation}
Given these observations, we can now analyze the employer's choices. 

\subsection*{Optimal solution for any ratio $\alpha\in(0,1)$}
%
The next theorem shows that for threshold policies, the  expected loss $l_{\alpha}(\pi)=l_{\alpha}(\theta)$ is minimized at $\theta^*_{p,\alpha}$ 
such that
$|\theta^*_{p,\alpha}-\tau/2|\leq \frac{\log (\frac{1}{p})+\log(\frac{1}{ \alpha})}{2\log(1+\frac{2\sigma}{1-\sigma})}$.
\begin{theorem}\label{thm:minErrorThresholdgeneralAlpha}
The loss function $l_{\alpha}(\theta)$ is quasi-convex and a
threshold of 
$\theta^*_{p,\alpha}=
\left\lceil\frac{1}{2}\left(\tau-\frac{\log(\frac{1}{p}-1)+\log(\frac{1}{\alpha}-1)}{\log(1+\frac{2\sigma}{1-\sigma})}\right)\right\rceil$ 
minimizes loss for any values of $\alpha,p,\sigma\in(0,1)$.
Namely,
\[
\theta^*_{p,\alpha}=\arg\min_{\theta}l_{\alpha}(\theta) =\left\lceil\frac{\tau}{2}-\frac{\log(\frac{1}{p}-1)+\log(\frac{1}{\alpha}-1)}{2\log(1+\frac{2\sigma}{1-\sigma})}\right\rceil.
\]
\end{theorem}

Next, we bound the number of tests required to guarantee that
the probability of classification error by the majority decision rule (i.e., $\theta=\ceil{\frac{\tau}{2}}$) 
does not exceed a specified quantity $\delta$.


\begin{theorem}\label{thm:majorityLowerBound}
For every $\delta,p,\alpha \in(0,1)$, performing 
$\tau=\Omega(\frac{\alpha+p-2p\alpha}{\sigma^2}\ln(\frac{1}{\delta}))$ 
tests per candidate 
and using majority as a decision rule (i.e., $\theta=\tau/2$)
guarantees $l_{\alpha}(\pi)\leq \delta$.
\end{theorem}

\subsubsection*{Equal cost for false positives and false negatives $(\alpha=\frac{1}{2})$}

Consider the simple loss consisting of the classification error rate 
(false positives and false negatives count equally), 
expressed via our loss function by setting $\alpha=\frac{1}{2}$. 
When skilled and unskilled candidates 
occur with equal frequency, i.e., $p=1/2$, 
we can derive that the majority decision rule 
minimizes the classification error 
for any number of tests.

\begin{corollary}\label{cor:majorityMinError}
Assume  $p=1/2$ and $\alpha=1/2$. 
For any number of tests $\tau$, 
the majority decision rule minimizes loss $l_\alpha$.
Namely,
$\arg\min_{\theta}l_{\frac{1}{2}}(\theta) =\ceil{\frac{1}{2}\tau}.$
In addition, for every $\delta\in(0,1)$, 
performing $\tau=\Omega(\frac{1}{\sigma^2}\ln(\frac{1}{\delta}))$ tests per candidate 
and using majority as a decision rule 
guarantees classification error with probability of at most $\delta$.
\end{corollary}


\paragraph*{FDR minimization with limited number of tests per hire for balanced groups}
Again, assuming balanced groups (i.e., $p=1/2$),
suppose that an employer would like 
to minimize the false discovery rate, 
subject to the constraint of lower bounding the hiring probability.
We can model this optimization problem 
by introducing a budget parameter $B>1$ 
to bound any predetermined (fixed) number of tests 
per hired candidate as follows:
\begin{equation}\label{optimization1}
\begin{aligned}
& \arg\min_{\pi}
& & \text{FDR}_{\pi}=\Pr[y_i=0|\Pr[\pi(\hat{y}_{i,1},\dots,\hat{y}_{i,\tau}) =1] \\
& \text{subject to}
& & \frac{\tau_\pi}{ \Pr[\pi(\hat{y}_{i,1},\dots,\hat{y}_{i,\tau}) =1]}\leq B
\end{aligned}
\end{equation}
where $\tau_\pi$ is the number of tests $\pi$ performs.
The following theorem shows that the optimal policy is a randomized threshold policy.
\begin{theorem}\label{thm:minFPwithBudget}
There exists a randomized threshold policy $\pi$ 
which is an optimal solution for (\ref{optimization1}).
\end{theorem}
\subsection{The Dynamic Policy (Adaptively-Allocated Tests)}
Recall that under a dynamic policy, 
the employer can decide after each test 
whether to accept, reject, or perform another test.
In general, dynamic policies are more efficient 
than those that must set a fixed number of tests. 
To build intuition,
consider a candidate that has passed $2$ out of $3$ tests.
As seen above, under an optimally-constructed fixed-test policy,
any candidate that fails a single test might be rejected.
\footnote{For example, if $B=18$ and $\eta=\frac{1}{3}$,
the lowest false discovery rate is achieved by $\tau=\theta=3$.}
However, the posterior probability 
that this candidate is in fact \emph{skilled}
may still be greater than that of a fresh candidate sampled from the pool.
Thus we can improve on the fixed-test policy by dynamically
allocating more tests to candidates 
until their posterior odds either 
dip below the prior odds or rise above the threshold for hiring.
%
The following theorem formalizes this notion
that it is better to administer more tests 
to a candidate that passed the majority of previous tests 
than to start afresh with a new candidate:


\begin{theorem}\label{thm:posteriorBetter}
For any $p,\sigma, \tau$, a candidate $i$ 
that passed $\theta>\frac{\tau}{2}$ out of $\tau$ tests 
is more likely to be a skilled
than a freshly-sampled candidate $i'$ 
for whom no test results are yet available, i.e.,
$\Pr[y_{i'}=+1]=p<\Pr[y_i=+1|S_{\tau}=\theta]$.
\end{theorem}
\begin{remark}
If $\theta<\frac{\tau}{2}$, the inequality would have been reversed.
\end{remark}


\paragraph*{The Greedy Policy}
% 
We now present a greedy algorithm 
that continues to test a candidate 
so long as the posterior probability that $y_i=+1$ 
is greater than $\epsilon'$ and smaller than $1-\epsilon$,
rejects a candidate whenever the posterior falls below $\epsilon'$
(absent fairness concerns, employers will set $\epsilon'=p$ for all groups), 
and accepts whenever the posterior rises above $1-\epsilon$. 
Given parameters $\epsilon,\epsilon'>0$, 
we show that the greedy policy solves the optimization problem 
of minimizing the mean number of tests under these constraints, i.e.,
\begin{equation*}
\begin{aligned}
& \underset{\tau}{\text{minimize}}
& & \E[\tau] \\
& \text{subject to}
& &\forall_i \pi(\hat{y}_{i,1},\dots,\hat{y}_{i,\tau}) =1 \text{ iff } \Pr[y_i=+1|\hat{y}_{i,1},\dots,\hat{y}_{i,\tau}]\geq 1-\epsilon\\
& & & \forall_i \pi(\hat{y}_{i,1},\dots,\hat{y}_{i,\tau}) =0 \text{ iff } \Pr[y_i=+1|\hat{y}_{i,1},\dots,\hat{y}_{i,\tau}]<\epsilon'
\end{aligned}
\end{equation*}
Our analysis of this policy builds upon the observation
that conditioned on a worker's skill,  
the posterior log-odds after each test
perform a one-dimensional random walk,
starting with the prior log-odds $\log(\frac{p}{1-p})$
and moving, after each test result, either left 
(upon a failed test) or right (upon a passed test).  
When (as in our model) the probability of a flip 
are equal for skilled and unskilled candidates, 
our random walk has a fixed step size.
Moreover, our random walk has \emph{absorbing barriers}
corresponding to (when $\epsilon'=p$) 
falling below the prior log odds (on the left)
and exceeding the hiring threshold (on the right).
Owing to the fixed step size and absorbing barriers,
our policy resembles the classic problem of Gambler's ruin,
in which a gambler wins or loses a unit of currency at each step,
and loses when crossing a threshold on the left (going bankrupt)
or on the right (bankrupting the opponent).
We formalize the random walk as follows where $X_j$
is the position on the walk at time $j$:
\begin{enumerate}
    \item $X_0$ is the prior log-odds of the candidate, i.e., $X_0=\log \frac{p}{1-p}$.
    \item After each test result, $\hat{y}_{i,j}$ is observed,
    $X_j=X_{j-1}+(2\hat{y}_{i,j}-1)\cdot
    \log \left(\frac{\Pr[\hat{y}_{i,j}=+1|y_i=+1]}{\Pr[\hat{y}_{i,j}=+1|y_i=0]}\right)$.
\end{enumerate}
Let $\pi_{Greedy}$ be the policy that accepts a candidate if $\Pr[y_i=+1|\hat{y}_{i,1},\dots,\hat{y}_{i,j}]\geq 1-\epsilon$, 
rejects if $\Pr[y_i=+1|\hat{y}_{i,1},\dots,\hat{y}_{i,j}]<\epsilon'$,
and otherwise conducts an additional test, i.e.,
\[
\pi_{Greedy}(\hat{y}_{i,1},\dots,\hat{y}_{i,j})=
 \begin{cases}
       0    &\quad\text{if } \Pr[y_i=+1|\hat{y}_{i,1},\dots,\hat{y}_{i,j}]<\epsilon' \\
       1    &\quad\text{if } \Pr[y_i=+1|\hat{y}_{i,1},\dots,\hat{y}_{i,j}]\geq 1-\epsilon\\
       \text{retest}    &\quad\text{else}\\
     \end{cases}.
\]
An employer will generally set
the lower absorbing barrier
to reject all candidates with posterior log odds less than $p$ 
since a fresh candidate from the pool is expected to be better.
However, when noise levels differ across groups,
we may prefer \emph{in the interest of fairness} 
to set $\epsilon'$ lower than $p$ for members of the noisier group, 
allowing us to equalize the frequency of false negatives across groups (see Section \ref{sec:twoGroups}).


\begin{lemma}\label{lemma:beta}
Let $\beta,\beta'\in \mathbb{R}$ be the parameters that satisfy $\frac{\beta}{\beta+1}=1-\epsilon$ and $\frac{\beta'}{\beta'+1}=\epsilon'$
(i.e., $\beta=\frac{1-\epsilon}{\epsilon}$and $\beta'=\frac{\epsilon'}{1-\epsilon'}$).
Then $X_{\tau}\geq \log \beta$ 
iff $\Pr[y_i=+1|\hat{y}_{i,1},\dots,\hat{y}_{i,\tau} ] \geq 1-\epsilon$ 
(iff the candidate is accepted) and $X_{\tau}< \log \beta'$ 
iff $\Pr[y_i=+1|\hat{y}_{i,1},\dots,\hat{y}_{i,\tau} ] < \epsilon'$ 
(iff the candidate is rejected).
\end{lemma}

\begin{corollary}\label{cor:piGreedy}
The policy $\pi_{Greedy}$ can be described as follows.
\[
\pi_{Greedy}(\hat{y}_{i,1},\dots,\hat{y}_{i,\tau})=
 \begin{cases}
       0     &\quad\text{if } X_{\tau} < \log \frac{\epsilon'}{1-\epsilon'} \\
       1    &\quad\text{if } X_{\tau}\geq \log(\frac{1-\epsilon}{\epsilon})\\
       \textnormal{retest}    &\quad\textnormal{else}\\
     \end{cases}
\]
\end{corollary} 

\noindent We use the following parameters in the next theorems: 
$$a=\left\lceil \frac{\log(\frac{(1-\epsilon)(1-\epsilon')(1+\sigma)}{\epsilon \epsilon'(1-\sigma)})}{\log(\frac{1+\sigma}{1-\sigma})}\right\rceil \gg \frac{1}{\sigma}
\text{\quad and \quad} z=
\left\lceil\frac{\log (\frac{p(1-\epsilon')(1+\sigma)}{\epsilon'(1-p)(1-\sigma)})}{\log(\frac{1+\sigma}{1-\sigma})}\right\rceil$$
\begin{theorem}[Expected number of tests per type]\label{thm:expectedNumOfTests}
The expected number of tests until a decision 
(namely accept or reject) for skilled candidates is $\E[\tau_{s}]=\frac{1}{\sigma}\left(a\cdot\frac{1-(\frac{1-\sigma}{1+\sigma})^{z}}{1-(\frac{1-\sigma}{1+\sigma})^{a}}-z\right)\approx\frac{2a}{1+\sigma}-\frac{z}{\sigma}$ and $\E[\tau_{u}]=\frac{1}{\sigma}\left(z-a\cdot\frac{1-(\frac{1+\sigma}{1-\sigma})^{z}}{1-(\frac{1+\sigma}{1-\sigma})^{a}}\right)\approx\frac{z}{\sigma}$ for unskilled candidates.
\end{theorem} 
For the probabilities of the candidates to be accepted or rejected, 
conditioned on their true skill level, 
we present the results in a form of confusion matrix in Table \ref{tab:confMat}. 


\begin{table}
\begin{center}
\caption{Confusion matrix for $\pi_{\text{greedy}}$
assuming $\epsilon\leq 1/4$ and $\epsilon'\leq p\leq 1/2$.
}
\small
\begin{tabular}{l|ll|ll}
\multicolumn{1}{c}{} &\multicolumn{2}{c}{\textbf{General $\epsilon'$}}&\multicolumn{2}{c}{\textbf{When $\epsilon' = p$}} \\ 
\multicolumn{1}{c}{} & 
\multicolumn{1}{c}{\textbf{Skilled}  ($y_i=+1$)} & 
\multicolumn{1}{c}{\textbf{Unskilled} ($y_i=0$)} &
\multicolumn{1}{c}{\textbf{Skilled} } & 
\multicolumn{1}{c}{\textbf{Unskilled} }
\\ 
\toprule
\textbf{accept} & $\text{TPR}=\Theta\left(1-\frac{\epsilon'}{p}(1-\sigma)\right)$ & $\text{FPR}=\Theta\left(\epsilon(p-\epsilon'+\epsilon'\sigma) \right)$  &$ \Theta(\sigma)$&$ \Theta(\epsilon p \sigma)$ \\
\textbf{reject} & $\text{FNR}=\Theta\left(\frac{\epsilon'}{p}(1-\sigma)\right)$   & $\text{TNR}=\Theta\left(1-\epsilon(p-\epsilon'+\epsilon'\sigma) \right)$ & $ \Theta(1-\sigma)$ &$\Theta(1-\epsilon p \sigma)$\\ 
\bottomrule
\end{tabular}
\label{tab:confMat}
\end{center}
\end{table}


\begin{theorem}\label{thm:expectedNumOfTestsUntilDecision}
The expected number of tests until deciding whether to accept or reject a candidate is $
\E[\tau|{\pi(y_{i,\tau})\in \{0,1\}}]
\approx \frac{ap}{\sigma}$, where 
$a\gg\frac{1}{\sigma}$.
\end{theorem}
